\section{Running It in Production}

\subsection{Docker Compose stacks}

Docker Compose describes several containers, their settings, networks and
volumes in one YAML file. The project has three:

\begin{description}
\item[\texttt{docker-compose.yml}] development services (Postgres, MinIO,
  registry); the API usually runs on the host for fast reloads.
\item[\texttt{sim/docker-compose.sim.yml}] an \emph{overlay} adding the
  containerized control plane and five simulated PCs, enabled with
  \verb|--profile sim|. Compose merges overlay files on top of the base file.
\item[\texttt{deploy/docker-compose.prod.yml}] the production stack:
  Postgres, MinIO, registry, API, portal and Caddy. Only Caddy publishes ports.
\end{description}

\paragraph{Required variables.} In the production file every secret is
written as \verb|${POSTGRES_PASSWORD:?set POSTGRES_PASSWORD}|. The
\texttt{:?} syntax makes Compose refuse to start, naming the variable, when it
is unset or empty. Forgetting a secret is a startup error, not a silent default.

\paragraph{Health checks and ordering.} Services declare health checks (for
example \texttt{pg\_isready} for Postgres), and dependants wait with
\texttt{condition: service\_healthy}. The API never starts before its database
accepts connections.

\subsection{Backups}

\texttt{deploy/backup.sh} saves three things:
\begin{itemize}
\item Postgres, with \texttt{pg\_dump -Fc} (a compressed, restorable dump);
\item the MinIO bucket, with \texttt{mc mirror};
\item the agent signing key. Losing it would mean every PC rejects every future
  update until reinstalled.
\end{itemize}

Registry images are deliberately left out: they can be rebuilt from their
requirements. A backup is only real once it has been restored, so the restore
script was tested: back up the development data, restore into a scratch
database and bucket, and compare. 306 jobs, 99 nodes and 481 objects matched.

\subsection{Alerts}

\texttt{control/app/alerts.py} computes the conditions an administrator must
hear about without watching a dashboard: a PC offline for ten minutes, low
disk, many lost leases in an hour, one PC failing repeatedly, a job queued too
long. New alerts are posted to a chat webhook, each with a one-hour cooldown so
a persistent problem does not spam.

\subsection{Runbooks}

\texttt{docs/runbooks.md} has twelve step-by-step procedures: stand up the
server, add a PC, canary an update, drain, quarantine, replace WSL, recover a
full disk, rotate credentials, back up and restore, clean up storage, daily
health check, and moving the server. A runbook is written so that someone
tired at 2 a.m. can follow it, and each step says how to check it worked.

\subsection{Deployment ladder}

\begin{enumerate}
\item Simulated lab passes the failure exam.
\item Control plane on an always-on server with production settings.
\item One canary PC: reboot without login, CUDA job, self-update, rollback.
\item The failure exam by hand on five real PCs.
\item Staged rollout 1, 5, 15, 30, 60 PCs, recording success rate and lease
  losses at each step.
\end{enumerate}

Each step limits how much can break at once: the same reasoning as rollout
rings, applied to the whole system.
